%% P19 --- Convexity and Jensen's inequality
%%
%% Every number in this program that the reader cannot do in their head comes
%% from code/p19_convexity_jensen.py and is pulled in with \val{}. The console
%% listing is a committed transcript written by that script.
%%
%% THE THESIS: convexity is a promise about the PROBLEM, not about the
%% algorithm. One basin means every local minimum is global, so the place an
%% optimiser stops is the answer rather than an answer. Deep learning breaks
%% that promise, and what survives is worth naming precisely.
%%
%% WHAT THIS PROGRAM IS OWED, read out of the files rather than remembered:
%%   F04  ALREADY WORKS THE PERPLEXITY DEMONSTRATION IN FULL -- both numbers,
%%        and which one a leaderboard should print -- and hands the general
%%        statement here BY NAME. So this program owes the INEQUALITY and its
%%        direction, not the demonstration. F04 also contrasts it with the
%%        macro/micro average: same smell, different mechanism, because there
%%        every step is linear and the WEIGHTS do the damage.
%%   F03  owns computing a perplexity from a mean loss and the base warning.
%%   F11  hands over "convexity, local against global minima" by name.
%%   F13  hands over "convexity and Jensen's inequality" by name, and owns the
%%        weighted average, which is the shape Jensen is about.
%%   P17  owns the second derivative and the bowl, so the convexity TEST is an
%%        object the reader already has doing a new job.
%%
%% WHAT IT LEAVES ALONE, checked against tools/programs.json:
%%   the optimisers that live with a non-convex surface           -> P20
%%   minibatch noise                                              -> P21
%%   constrained optimisation and the multiplier                  -> P22
%%   what a perplexity MEANS (the effective number of choices)    -> P29
%%   KL, Jensen-Shannon, the asymmetry                            -> P30
%%
%% THE HEADLINE IS THAT THE ERROR'S SIZE IS SET BY THE EVALUATION SET. The
%% ratio of the wrong average to the right one is exp(Var/2) for Gaussian
%% losses: 1.00 on a homogeneous corpus and 7.42 at a spread of two nats. So a
%% harness carrying the bug passes every test on a toy set and misreports on a
%% real one, which is P02's "safe for inputs you have not tried" arriving in
%% published evaluation code.
%%
%% Twin: programs/pl/P19-convexity-jensen.tex. Frame for frame, macro for
%% macro, number for number; tools/parity.py compares them.

\program[Convexity and Jensen]{Convexity and Jensen's inequality}
\label{prog:P19}
\index{convexity}
\index{Jensen's inequality}

\begin{outcomes}
\outcome{1--8}{state the chord definition of convexity and apply it as a test
  rather than as a picture}
\outcome{9--14}{say what convexity promises, and say precisely what the
  promise is about}
\outcome{15--18}{state Jensen's inequality and get its direction right for a
  convex and for a concave function}
\outcome{19--26}{say why an average of perplexities is a different quantity,
  and what decides how different}
\outcome{27--30}{recognise the same inequality as the ELBO and as a variance
  being non-negative}
\outcome{31--38}{say what is left when the promise is broken, without either
  overclaiming or calling it hopeless}
\end{outcomes}

\begin{quiz}
\item What does convexity promise about a minimum? \teachesat{9--10}
  \answerto{That every local minimum is global. It is a promise about the
  problem rather than about the algorithm: it says the place an optimiser
  stops is \emph{the} answer, and says nothing about how long it takes to get
  there.}
\item A curved function is applied to an average. Is that the same as
  averaging the function's values? \teachesat{17--18}
  \answerto{No, and Jensen's inequality says which way round: for a convex
  function the average of the values is the larger. They agree only when
  every value is the same.}
\item A leaderboard averages the per-document perplexities of a model. What
  is it reporting? \teachesat{19--20}
  \answerto{A different quantity from the perplexity, and a larger one. The
  perplexity is the exponential of the mean loss; the mean of the
  exponentials sits above it by a factor that grows with how varied the
  documents are.}
\item Is $\ln$ convex or concave, and what does that do to the inequality?
  \teachesat{27--28}
  \answerto{Concave, so the inequality turns round: the logarithm of a mean
  is at least the mean of the logarithms. That one line is the ELBO.}
\item Deep learning is not convex. What is lost, and what is not?
  \teachesat{33--34}
  \answerto{What is lost is the guarantee that where you stopped is the best
  there is. What is not lost is everything else: the gradient still points
  downhill, the step-size bound still holds, and the counting argument says
  most stationary points are saddles rather than the bad minima people
  picture.}
\end{quiz}

% ============================================================
\section{A definition you can check}
\label{sec:P19-chord}
\index{convexity!the chord definition}

\begin{fr}
Programs \ref{prog:F11}, \ref{prog:P15} and \ref{prog:P17} all leant on the
same picture without ever naming it: a bowl, with a bottom, and a walk that
goes down. Program~\ref{prog:F11} said in as many words that
\enquote{convexity, local against global minima} was owed here.

Start with what the word has to mean before any of that pays off. Draw a
curve that is bowl-shaped, pick two points on it, and join them with a
straight line. Where does that line sit?
\dotline
\nextframe
\end{fr}

\begin{fr}
\begin{ansblock}
Above the curve, or on it \dash{} never below.
\end{ansblock}

That line is the \textbf{chord}, and what you just said is the definition. A
function is \textbf{convex} when, for every pair of points and every position
between them,
\[ f\bigl(t a + (1-t)b\bigr) \;\le\; t\,f(a) + (1-t)\,f(b) \]
for $t$ between $0$ and $1$. The left is the curve, the right is the chord.

\mermaidfig{p19-chord-above}{Two points, and the line between them.}{the chord
test}
\end{fr}

\begin{fr}
That inequality is worth more than the picture it came from, and the reason is
a practical one rather than an aesthetic one.

What can you do with the inequality that you cannot do with the bowl?
\nextframe
\end{fr}

\begin{fr}
\begin{ansblock}
Check it. And check it in a number of dimensions nobody can draw.
\end{ansblock}

A picture of a bowl is an argument about two dimensions and a hope about the
rest. The inequality is arithmetic: it takes two points and a number between
$0$ and $1$, and it either holds or it does not.

The script does exactly that for $\val{p19.chord.fns}$ functions over
$\val{p19.chord.tests}$ triples each \dash{} $x^{2}$, $e^{x}$ and
$\lvert x \rvert$ \dash{} and every one holds.

\begin{note}
$\lvert x \rvert$ is in that list on purpose. It is convex and it has a corner,
so it is not differentiable at the origin. \textbf{Convexity is not a
smoothness condition}, which matters because $\relu$ is exactly this shape on
one side and the loss surfaces this book is about are full of it.
\end{note}
\end{fr}

\begin{fr}
A test that cannot fail is not a test, which is Program~\ref{prog:P05}'s rule
about finding nothing.

So the script also runs it on $x^{4} - 3x^{2} + \frac{x}{2}$, which is not
convex. What should the sweep report?
\dotline
\nextframe
\end{fr}

\begin{fr}
\begin{ansblock}
A chord that dips \emph{below} the curve somewhere \dash{} here by
$\val{p19.chord.counterexample}$ at the worst pair.
\end{ansblock}

Which is not a rounding failure or a near miss. The definition is violated by
a margin, because that quartic has two dips with a rise between them, and a
chord drawn across the rise passes underneath it.

That function is going to do the rest of this program's work, so it is worth
looking at it once: it is the smallest thing that behaves the way a loss
surface does.
\end{fr}

\begin{fr}
There is a second test, and it is an object you already have.

Program~\ref{prog:P17} built the second derivative and used it to classify a
stationary point. Use it here instead: what must be true of $f''$ for the
chord to lie above the curve everywhere?
\dotline
\nextframe
\end{fr}

\begin{fr}
\begin{ansblock}
$f'' \ge 0$ everywhere \dash{} the curve never bends downwards.
\end{ansblock}

And in more than one dimension it is the same statement about
Program~\ref{prog:P17}'s matrix: the Hessian is positive semi-definite at
every point, which by Program~\ref{prog:P10} means every eigenvalue is
non-negative in every direction, everywhere.

Measured on the quartic: it curves downwards at $\val{p19.curv.negative}$ of
$\val{p19.curv.tested}$ grid points. One would be enough.

\begin{note}
Note the quantifier, because it is where the two tests differ in practice.
Program~\ref{prog:P17} asked about the Hessian \emph{at a point} and got a
classification of that point. Convexity asks about the Hessian \emph{at every
point} and gets a statement about the whole function. Same object, and a
much stronger demand.
\end{note}
\end{fr}

% ============================================================
\section{What convexity promises}
\label{sec:P19-promise}
\index{convexity!and local minima}

\begin{fr}
Now the payoff, and it is one sentence.

You are optimising. Your walk stops somewhere, because the ground has stopped
going down. On a convex surface, what do you know about where you are?
\nextframe
\end{fr}

\begin{fr}
\begin{ansblock}
That it is the best there is. Every local minimum of a convex function is
global.
\end{ansblock}

The argument is the chord again, and it is two lines. Suppose $a$ is a local
minimum and some $b$ is lower. Every point on the chord from $a$ to $b$ is at
most the chord's height, which slopes down from $f(a)$ towards $f(b)$ \dash{}
so points arbitrarily close to $a$ are strictly below $f(a)$, and $a$ was not
a local minimum after all.

\mermaidfig{p19-one-basin}{Where the walk stopped, and what that is worth.}{one
basin}
\end{fr}

\begin{fr}
The measurement is a count, and it is deliberately unimpressive.

$x^{2}$ has $\val{p19.basins.convex}$ local minimum on the interval the script
sweeps. The quartic has $\val{p19.basins.wiggle}$. Say what the second number
costs you.
\dotline
\nextframe
\end{fr}

\begin{fr}
\begin{ansblock}
That where you start decides what you get. The two minima are worth
$\val{p19.wiggle.lo}$ and $\val{p19.wiggle.hi}$ \dash{} a gap of
$\val{p19.wiggle.gap}$ \dash{} and a walk finds whichever one it was nearest.
\end{ansblock}

\textbf{So convexity is not a statement about difficulty. It is a statement
about ambiguity.} A convex problem can be enormous and slow; what it cannot be
is a problem where two honest runs disagree about the answer and both are
right.

\begin{trapbox}
\textbf{\enquote{Convex means easy to optimise.}} It means the answer is
unique, which is a different property and neither implies the other.

Convex problems can be slow: Program~\ref{prog:P17} measured a bowl whose
condition number forced $\val{p17.steps.99}$ steps to close $99$ per cent of
one direction's gap, and that bowl is as convex as anything gets. And
non-convex problems can be fast, which is the entire empirical history of
deep learning.

The promise is about \emph{which} point you end at, not about how long it
takes to arrive.
\end{trapbox}
\end{fr}

\begin{fr}
One more thing the promise is not, and it is the one that gets misquoted.

Convexity is a property of the \emph{function}. Say what it therefore tells
you about gradient descent.
\nextframe
\end{fr}

\begin{fr}
\begin{ansblock}
Nothing directly. It is a promise about the problem, not about the algorithm.
\end{ansblock}

An optimiser can still diverge on a convex function \dash{}
Program~\ref{prog:P17}'s bound is $\eta < 2/\lambda_{\max}$ whether or not
the surface is convex \dash{} and it can still crawl. What convexity buys is
that \emph{if} the walk converges to a stationary point, that point is the
answer.

Separating those two is the whole of this section, and it is why the sentence
\enquote{the loss is non-convex so training is unreliable} contains two claims
of which only one is about convexity.
\end{fr}

% ============================================================
\section{Jensen's inequality}
\label{sec:P19-jensen}
\index{Jensen's inequality!direction of}

\begin{fr}
The chord definition has a second reading, and it is the one this program was
written for.

Read $t\,f(a) + (1-t)\,f(b)$ as a \emph{weighted average of two values of
$f$}, and $t a + (1-t) b$ as the same weighted average of $a$ and $b$. Now say
what the definition is comparing.
\dotline
\nextframe
\end{fr}

\begin{fr}
\begin{ansblock}
$f$ of an average, against the average of $f$. And it says the first is at
most the second.
\end{ansblock}

Which is Program~\ref{prog:F13}'s weighted average with a function wrapped
round it, and the general form is \textbf{Jensen's inequality}: for a convex
$f$ and any weights that are non-negative and sum to one,
\[ f\Bigl(\sum_i w_i x_i\Bigr) \;\le\; \sum_i w_i\, f(x_i) \]

Two points was the definition; $n$ points is the same statement applied
repeatedly, and it needs nothing new.
\end{fr}

\begin{fr}
The direction is the part everybody gets wrong, and it is worth pinning down
rather than remembering.

You have a convex $f$. Which is bigger \dash{} averaging first, or applying
$f$ first? And when are they equal?
\nextframe
\end{fr}

\begin{fr}
\begin{ansblock}
Applying $f$ first is bigger. They are equal exactly when every $x_i$ is the
same.
\end{ansblock}

The equality case is the useful half. Jensen's gap is not a fixed penalty: it
is zero when the inputs are identical and grows as they spread out, which is
what makes it a fact about \emph{your data} rather than about the function.

\begin{note}
For a \textbf{concave} $f$ every inequality here turns round, because a
concave function is $-1$ times a convex one and multiplying by a negative
number reverses an inequality. So $\ln$ of a mean is at least the mean of the
$\ln$s, which \S\ref{sec:P19-elbo} uses.
\end{note}
\end{fr}

% ============================================================
\section{Why you cannot average perplexities}
\label{sec:P19-perplexity}
\index{perplexity!averaging}

\begin{fr}
Program~\ref{prog:F04} set this up and handed the general statement here by
name. It worked a four-example cross-entropy, exponentiated the mean loss to
get $\val{f04.ce.ppl}$, averaged the four per-example perplexities to get
$\val{f04.ce.ppl.naive}$, and said the first is what a leaderboard should
print.

You now have the reason. Say which of those two numbers must be the larger,
and why it is not a coincidence of that example.
\dotline
\nextframe
\end{fr}

\begin{fr}
\begin{ansblock}
The average of the perplexities, always \dash{} because $e^{x}$ is convex and
Jensen says so for every set of losses.
\end{ansblock}

So the two numbers are \textbf{ordered}, not merely different, and the wrong
one is always the larger. A harness with this bug always flatters nothing and
always reports a model as worse than it is.

The script asserts that ordering rather than the figures, over
$\val{p19.ppl.tokens}$ tokens at four different spreads, and it reads
Program~\ref{prog:F04}'s two committed values and checks they are ordered the
same way \dash{} so the two programs cannot come apart about which number is
which.
\end{fr}

\begin{fr}
\begin{trapbox}
\textbf{\enquote{Perplexity is the average of the per-token perplexities.}}
It is the exponential of the average loss, which is not the same number, and
the difference is not small.

The reason the belief survives is that it is \emph{almost} true, in a precise
sense: the two agree exactly when every token costs the same. So an
implementation carrying this error is correct on a uniform test fixture, and
that is usually what it was tested on.
\end{trapbox}
\end{fr}

\begin{fr}
Which raises the question the rest of this section is for.

The gap is zero when every loss is the same, and it grows as they spread out.
So what property of your evaluation set decides how wrong the wrong number is?
\nextframe
\end{fr}

\begin{fr}
\begin{ansblock}
How varied it is \dash{} the spread of the per-token losses, and nothing
else about the model.
\end{ansblock}

That is the finding, and it is sharper than \enquote{the two numbers differ}.
Same mean loss, $\val{p19.ppl.meanloss}$, four different spreads:

\begin{center}
\begin{tabular}{lrrr}
\toprule
spread & right & wrong & wrong / right \\
\midrule
$0$ & $\val{p19.ppl.right.00}$ & $\val{p19.ppl.wrong.00}$ & $\val{p19.ppl.ratio.00}$ \\
$\num{0.5}$ & $\val{p19.ppl.right.05}$ & $\val{p19.ppl.wrong.05}$ & $\val{p19.ppl.ratio.05}$ \\
$\num{1.0}$ & $\val{p19.ppl.right.10}$ & $\val{p19.ppl.wrong.10}$ & $\val{p19.ppl.ratio.10}$ \\
$\num{2.0}$ & $\val{p19.ppl.right.20}$ & $\val{p19.ppl.wrong.20}$ & $\val{p19.ppl.ratio.20}$ \\
\bottomrule
\end{tabular}
\end{center}

The model is identical down the column: the right number does not move. The
reported one goes from correct to $\val{p19.ppl.ratio.20}$ times too large.

\mermaidfig{p19-order-of-two}{One column, two possible quantities.}{order of
two steps}
\end{fr}

\begin{fr}
\textbf{So the size of the error is a property of the evaluation set and not
of the code}, which is the part that makes this worth a section.

A harness with this bug is exactly right on a homogeneous fixture, mildly
wrong on a clean benchmark, and wrong by a factor on a diverse corpus \dash{}
which is the corpus anybody actually cares about. It passes its tests and
misreports in production, and that is Program~\ref{prog:P02}'s
\enquote{numerically stable} in a new place: not more accurate, but safe on
inputs you have not tried.

\begin{note}
There is a closed form for the ratio when the losses are near-Gaussian
\dash{} it is $e^{\sigma^{2}/2}$ \dash{} and the four rows above agree with
it to within $\val{p19.ppl.pred.bound}$ per cent. The script \emph{reports}
that and does not assert it, because it is exact for a log-normal and only
second-order otherwise, and asserting a law that holds in one case is how a
plausible number gets into a book.
\end{note}
\end{fr}

\begin{fr}
Here it is in four lines, on two sets of losses with the same mean.

\transcript{p19-two-averages}

Both columns of the first pair agree. Say what changed in the second, and what
did not.
\dotline
\nextframe
\end{fr}

\begin{fr}
\begin{ansblock}
The mean loss did not change, so the perplexity did not. The spread changed,
and the average of the perplexities nearly quadrupled.
\end{ansblock}

Four numbers is enough to show it because the mechanism is not statistical.
It is the shape of $e^{x}$: the high-loss tokens are exponentiated before they
are averaged, so they dominate a sum they would only have contributed to.

\begin{note}
Keep this apart from the failure Program~\ref{prog:F04} sets beside it. There,
a macro average and a micro average disagree because the groups have different
sizes and every step is \emph{linear}; the damage is done by the weights.
Here every weight is equal and the damage is done by the curvature. Same
smell, different mechanism, and they do not have the same fix.
\end{note}
\end{fr}

% ============================================================
\section{The same inequality, once more}
\label{sec:P19-elbo}
\index{Jensen's inequality!and the ELBO}

\begin{fr}
Jensen turns up in one more place a reader of this field will meet, and the
point of this section is recognition rather than derivation.

$\ln$ is concave. Write down what Jensen therefore says about
$\ln$ of an average.
\nextframe
\end{fr}

\begin{fr}
\begin{ansblock}
\[ \ln\Bigl(\sum_i w_i x_i\Bigr) \;\ge\; \sum_i w_i \ln x_i \]
\end{ansblock}

Checked over $\val{p19.elbo.trials}$ random positive samples, and on the
worked set in the script the gap is $\val{p19.elbo.gap}$.

\textbf{That one line is the evidence lower bound.} Written with an
expectation rather than a weighted sum it reads
$\ln \Ex[x] \ge \Ex[\ln x]$, and every variational method in the field is
built on choosing the $x$ so that the right-hand side is something you can
compute and the left-hand side is something you cannot.

\begin{rigourbox}
This program names the bound and does not use it. What the $x$ is, why the
bound is tight when it is, and what is being approximated are a course on
variational inference, and this book does not undertake one. What is here is
enough to recognise the move when a paper makes it in one line and expects
you to follow.
\end{rigourbox}
\end{fr}

\begin{fr}
And one that costs nothing to notice, because it is the same inequality with
$f(x) = x^{2}$.

Jensen says the mean of the squares is at least the square of the mean. Say
what the difference between them is called.
\dotline
\nextframe
\end{fr}

\begin{fr}
\begin{ansblock}
The variance.
\end{ansblock}

$\Ex[x^{2}] - (\Ex[x])^{2}$ is the variance, and Jensen's inequality for
$x^{2}$ is exactly the statement that it cannot be negative. Which is not a
new fact; it is a fact you already had, arriving from a direction that makes
it obvious.

\textbf{That is the case for knowing Jensen by name.} It is one inequality
that turns up as the perplexity error, as the ELBO, and as a variance being
non-negative, and recognising it is cheaper than meeting all three
separately.
\end{fr}

% ============================================================
\section{When the promise is broken}
\label{sec:P19-nonconvex}
\index{convexity!what survives without it}

\begin{fr}
The loss surface of a neural network is not convex, and it is not close.

Program~\ref{prog:P06} says why in one line, and it is worth recovering.
A network is a composition, and a composition of convex functions is not
generally convex. Give the shortest reason.
\nextframe
\end{fr}

\begin{fr}
\begin{ansblock}
Because the loss is a function of the \emph{weights}, and the weights of
different layers multiply each other.
\end{ansblock}

Even two layers with no activation give a loss containing $w_2 w_1$, and a
product of two variables is a saddle: convex along one diagonal, concave along
the other. The non-convexity is there before any activation function, before
any depth, and it comes from the thing that makes a network a network.

So the promise is gone. The rest of this section is about what is not.
\end{fr}

\begin{fr}
Start with what is actually lost, stated narrowly. What exactly stops being
true?
\dotline
\nextframe
\end{fr}

\begin{fr}
\begin{ansblock}
That where the walk stopped is the best there is. Nothing else.
\end{ansblock}

Everything this book has derived about optimisation was derived without
convexity anywhere in the argument: Program~\ref{prog:P15}'s steepest
direction is a fact about a dot product, Program~\ref{prog:P17}'s
$\eta < 2/\lambda_{\max}$ is a fact about a geometric sequence, and
Program~\ref{prog:P18}'s gradients are identities. None of them asked whether
the surface had one basin.

\textbf{So \enquote{the loss is non-convex} does not license pessimism about
any of them.} It licenses exactly one doubt, about one claim.
\end{fr}

\begin{fr}
And now the honest part, because the picture the phrase conjures is wrong in a
specific way.

Most people hear \enquote{non-convex} and picture the quartic from
\S\ref{sec:P19-chord} in a billion dimensions: a landscape of bad local
minima, and training as the luck of which one you land in.
Program~\ref{prog:P17} has an argument that says this is not the shape of the
problem. What was it?
\dotline
\nextframe
\end{fr}

\begin{fr}
\begin{ansblock}
The counting argument: a minimum needs \emph{every} eigenvalue sign to be
positive, and a saddle needs only one to differ, so in high dimension a
stationary point is overwhelmingly likely to be a saddle.
\end{ansblock}

A saddle has a downhill direction. So the thing that makes a run slow is
usually a place with a way out that the gradient is too small to find
quickly, rather than a basin with no way out at all \dash{} which is a
different problem with different remedies.

\begin{warning}
\textbf{This book has not measured the loss surface of any trained network,
and the paragraph above is arithmetic about signs rather than a finding about
models.} What it establishes is that the two-dimensional picture does not
transfer, not what the real geometry is.

What the literature reports \dash{} that the minima reached in practice tend
to be close in value \dash{} is an empirical claim about particular
architectures and particular data, and it is not settled here. The filter
Program~\ref{prog:P17} supplies applies to it: ask what quantity is being
compared, and whether it survives a reparameterisation.
\end{warning}
\end{fr}

\begin{fr}
One last question, because it is the one worth taking away.

Convexity is gone and cannot be recovered. So what is the point of a program
about it?
\nextframe
\end{fr}

\begin{fr}
\begin{ansblock}
That the inequality survives the promise. Jensen holds for every convex
function whether or not the thing you are optimising is one.
\end{ansblock}

\textbf{The promise is about a surface and the inequality is about an
average}, and the second is used far more often than the first. Every time
this field averages something and then transforms it \dash{} or transforms
and then averages \dash{} the order matters, and Jensen says which way.

The perplexity error in \S\ref{sec:P19-perplexity} is that mistake made in
public. The ELBO is that inequality used on purpose. Both are the same three
lines of \S\ref{sec:P19-jensen}, and neither needs the loss surface to be
anything at all.
\end{fr}

% ============================================================
\begin{summarybox}
\sumitem{1--2}{\result{\textbf{A function is convex when the chord never dips below
  the curve}, for every pair of points and every position between them. That
  is the definition, not a picture of it}.}
\sumitem{3--4}{\result{\textbf{The inequality can be checked where the picture cannot},
  which is why it is the useful form: $\val{p19.chord.fns}$ functions over
  $\val{p19.chord.tests}$ triples each. $\lvert x\rvert$ is convex and has a
  corner, so \textbf{convexity is not smoothness}}.}
\sumitem{5--6}{A test that cannot fail is not a test: the same sweep on a
  quartic fails by $\val{p19.chord.counterexample}$.}
\sumitem{7--8}{\result{\textbf{The second test is Program~\ref{prog:P17}'s object}:
  $f'' \ge 0$ everywhere, or the Hessian positive semi-definite everywhere.
  Note the quantifier \dash{} P17 asked at a point, this asks at every point}.}
\sumitem{9--10}{\result{\textbf{Every local minimum of a convex function is global},
  by the chord in two lines. That is what the promise is}.}
\sumitem{11--12}{\result{$\val{p19.basins.convex}$ basin against
  $\val{p19.basins.wiggle}$, worth $\val{p19.wiggle.lo}$ and
  $\val{p19.wiggle.hi}$. \textbf{Convexity is about ambiguity, not
  difficulty}: a convex problem can be slow and a non-convex one fast}.}
\sumitem{13--14}{\result{\textbf{It is a promise about the problem, not the
  algorithm.} An optimiser can still diverge or crawl on a convex function;
  what convexity buys is that a stationary point it reaches is the answer}.}
\sumitem{15--16}{\result{\textbf{Jensen's inequality} is the chord read as an average:
  $f$ of a weighted mean is at most the weighted mean of $f$}.}
\sumitem{17--18}{\result{\textbf{Applying the convex function first gives the larger
  number}, and the two are equal exactly when every input is the same. For a
  concave function every inequality turns round}.}
\sumitem{19--20}{\result{\textbf{An average of perplexities is always the larger of
  the two}, because $e^{x}$ is convex \dash{} so the two numbers are ordered
  and not merely different, and Program~\ref{prog:F04}'s pair is gated
  against that ordering}.}
\sumitem{21}{\result{\textbf{The belief survives because it is almost true}: the two
  agree exactly when every token costs the same, which is what a test fixture
  looks like}.}
\sumitem{22--24}{\result{\textbf{The size of the error is set by the spread of the
  losses}, not by the model: $\val{p19.ppl.ratio.00}$ at zero spread and
  $\val{p19.ppl.ratio.20}$ at $\num{2.0}$ nats, with the right number
  unmoved. A harness with the bug passes on a fixture and misreports on a
  real corpus}.}
\sumitem{25--26}{\result{Keep it apart from the macro/micro average
  Program~\ref{prog:F04} sets beside it: \textbf{there the weights do the
  damage and here the curvature does}, and the two do not have the same fix}.}
\sumitem{27--28}{\result{\textbf{$\ln$ is concave, so the inequality turns round}, and
  $\ln\Ex[x] \ge \Ex[\ln x]$ is the ELBO in one line.} Named here for
  recognition; the method is not this book's.}
\sumitem{29--30}{\result{\textbf{Jensen for $x^{2}$ is the variance being
  non-negative}, which is why the inequality is worth knowing by name: one
  statement, three places}.}
\sumitem{31--32}{\result{\textbf{A network's loss is non-convex before any activation
  or any depth}, because the weights of different layers multiply and a
  product of two variables is a saddle}.}
\sumitem{33--38}{\result{\textbf{What is lost is one claim and nothing else.} The
  steepest direction, the step-size bound and the gradient identities were all
  derived without convexity. \textbf{And the inequality survives the promise},
  which is why this program is worth its frames}.}
\end{summarybox}

\canyou

\begin{testexercises}
\item A colleague says their loss is convex, so the optimiser will converge
  quickly. Give the two things wrong with that sentence.
  \answerto{Convexity says nothing about speed \dash{} a convex bowl with a
  large condition number is slow, which Program~\ref{prog:P17} measured. And
  it says nothing about whether the optimiser converges at all: a step size
  above $2/\lambda_{\max}$ diverges on a convex function as readily as on any
  other. What it promises is that a stationary point reached is the global
  minimum.}
\item Two harnesses report perplexities of $\val{p19.ppl.right.10}$ and
  $\val{p19.ppl.wrong.10}$ for the same model on the same corpus. Which is
  more likely to be right, and what would settle it?
  \answerto{The smaller, because Jensen puts the average of perplexities
  above the exponential of the mean loss and never below. What settles it is
  running both on a corpus where every token costs the same: there the two
  agree, and a harness that still reports two numbers has a different bug.}
\item Is $f(x) = x^{3}$ convex on the whole line? Give the test you used.
  \answerto{No. $f'' = 6x$, which is negative for every $x < 0$, so the
  second-derivative test fails on the whole left half. It is convex on
  $x \ge 0$ \dash{} convexity is a property of a function \emph{and a
  domain}, and dropping the domain is how the word gets misused.}
\item Your evaluation set is one document type today and five tomorrow. What
  happens to a perplexity-averaging bug, and why is that the worst possible
  behaviour for a bug?
  \answerto{It gets worse, because the spread of the per-token losses grows
  and the ratio grows with it. It is the worst behaviour because the bug is
  quietest exactly when the test set is simplest: it passes on the fixture,
  survives review, and misreports on the corpus that was the point of the
  exercise.}
\item Give an example of a convex function with no minimum at all, and say
  what that shows about the promise.
  \answerto{$e^{x}$, or any straight line with non-zero slope: convex
  everywhere, decreasing forever, no minimum. It shows the promise is
  conditional \dash{} \emph{if} a local minimum exists it is global. Convexity
  does not assert that one exists, which is a separate question about the
  function and its domain.}
\end{testexercises}

\begin{furtherproblems}
\item Show that the sum of two convex functions is convex, and that the
  maximum of two convex functions is convex.
  \answerto{For the sum, add the two chord inequalities: each says the chord
  is at least the curve, and adding preserves that. For the maximum, at any
  point the max is one of the two, and that one's chord lies above it while
  the max of the two chords lies above that chord. Both are worth having
  because they are how convex objectives get built: a sum of per-example
  losses and a hinge are exactly these two operations.}
\item A model's per-token losses on a corpus are all equal. What does the
  perplexity-averaging bug report, and what does that say about testing for
  it?
  \answerto{Exactly the right answer, because Jensen's gap is zero when every
  input is the same. Which says a unit test on uniform data cannot detect the
  bug at all: the test has to use losses that differ, and it should assert the
  \emph{ordering} rather than a figure, since the size depends on the spread.}
\item Program~\ref{prog:F04} shows a macro average and a micro average
  disagreeing. Could Jensen's inequality tell you which is larger?
  \answerto{No, and that is the point of keeping them apart. Jensen needs a
  convex function applied to an average; there every operation is linear and
  the disagreement comes from unequal group sizes, so either can be the
  larger depending on which groups are big. The two failures look alike in a
  report and have different causes and different fixes.}
\item The arithmetic mean is at least the geometric mean. Derive it from
  Jensen.
  \answerto{Apply Jensen to the concave $\ln$: $\ln$ of the arithmetic mean is
  at least the mean of the $\ln$s, which is the $\ln$ of the geometric mean.
  Exponentiating both sides preserves the direction, so the arithmetic mean is
  at least the geometric mean, with equality exactly when every value is the
  same. That is the same inequality as this program's perplexity result read
  in the other direction.}
\item Why does this program prove nothing about whether the minima a real
  network reaches are close in value?
  \answerto{Because that is an empirical claim about particular architectures
  and data, and nothing here measures a trained network. What the program
  supplies is arithmetic about signs \dash{} that a stationary point in high
  dimension is overwhelmingly likely to be a saddle \dash{} which says the
  two-dimensional picture of bad basins does not transfer, and does not say
  what the real geometry is. Program~\ref{prog:P17}'s filter applies: ask what
  quantity is compared and whether it survives a reparameterisation.}
\end{furtherproblems}
